% Part VII — Architecture and Ecosystem (Slides 161–180)
\providecommand{\jugmark}[1]{\textbf{\color{jugeoBlue}#1}}

% ---------------------------------------------------------------------------
% Slide 161: Section divider
% ---------------------------------------------------------------------------
\begin{frame}
\frametitle{}
\vfill
\begin{center}
  {\Huge\color{jugeoBlue}\textbf{Part VII}}\\[12pt]
  {\LARGE\color{jugeoDark}Under the Hood --- Architecture}\\[20pt]
  \begin{tikzpicture}
    \draw[jugeoBlue, line width=2pt] (-3,0) -- (3,0);
  \end{tikzpicture}\\[12pt]
  {\large\color{gray}\itshape 960 modules, 10 layers, one dependency DAG}
\end{center}
\vfill
\end{frame}

% ---------------------------------------------------------------------------
% Slide 162: Scale of JuGeo
% ---------------------------------------------------------------------------
\begin{frame}
\frametitle{Scale of JuGeo}
\begin{center}
\begin{tikzpicture}[
  bignum/.style={font=\fontsize{48}{56}\selectfont\bfseries\color{jugeoBlue}},
  caption/.style={font=\large\color{jugeoDark}}
]
  \node[bignum] (n) {$\sim$960};
  \node[caption, below=4pt of n] {Python modules};
\end{tikzpicture}
\end{center}
\vfill
\begin{itemize}
  \item Organized in a strict \jugmark{dependency DAG} (directed acyclic graph)
  \item Each layer builds only on the layers below it
  \item No circular imports, no upward calls
  \item Like the layers of an operating system
\end{itemize}
\end{frame}

% ---------------------------------------------------------------------------
% Slide 163: The Architecture DAG — Overview
% ---------------------------------------------------------------------------
\begin{frame}
\frametitle{The Architecture DAG --- Overview}
\begin{center}
\begin{tikzpicture}[
  layer/.style={draw=jugeoBlue, fill=jugeoLight, rounded corners,
                minimum width=6cm, minimum height=7mm,
                font=\small\bfseries, text=jugeoDark, line width=0.8pt},
  arr/.style={-{Stealth[length=3mm]}, jugeoBlue!60, line width=0.7pt}
]
  \node[layer, fill=jugeoRed!15]   (pm) at (0,5.4) {problem\_modes};
  \node[layer, fill=jugeoOrange!15] (or) at (0,4.5) {orchestration / ideation};
  \node[layer, fill=jugeoOrange!10] (ge) at (0,3.6) {generation};
  \node[layer]                      (pr) at (0,2.7) {python\_runtime};
  \node[layer]                      (en) at (0,1.8) {encodings};
  \node[layer]                      (so) at (0,0.9) {solver};
  \node[layer]                      (ev) at (0,0.0) {evidence};
  \node[layer]                      (ju) at (0,-0.9){judgments};
  \node[layer]                      (go) at (0,-1.8){geometry};
  \node[layer, fill=jugeoGreen!15]  (ke) at (0,-2.7){kernel};

  \foreach \a/\b in {pm/or, or/ge, ge/pr, pr/en, en/so, so/ev, ev/ju, ju/go, go/ke} {
    \draw[arr] (\a) -- (\b);
  }

  \node[font=\scriptsize\color{gray}, right=6pt of pm] {Layer 10};
  \node[font=\scriptsize\color{gray}, right=6pt of ke] {Layer 1};
\end{tikzpicture}
\end{center}
\end{frame}

% ---------------------------------------------------------------------------
% Slide 164: Layer 1 — Kernel
% ---------------------------------------------------------------------------
\begin{frame}
\frametitle{Layer 1 --- Kernel}
\begin{columns}[T]
\begin{column}{0.55\textwidth}
\begin{block}{What It Does}
  \begin{itemize}
    \item Core services and lifecycle management
    \item Authority and trust anchoring
    \item The \jugmark{trusted computing base}
  \end{itemize}
\end{block}
\vspace{6pt}
\begin{exampleblock}{Analogy}
  Like an OS kernel --- \textbf{small}, \textbf{trusted},\\
  and everything else builds on it.
\end{exampleblock}
\end{column}
\begin{column}{0.42\textwidth}
\begin{center}
\begin{tikzpicture}[
  ring/.style={draw=jugeoBlue, line width=1pt}
]
  \fill[jugeoGreen!20] (0,0) circle (1.8);
  \fill[jugeoLight] (0,0) circle (1.2);
  \fill[jugeoGreen!40] (0,0) circle (0.6);
  \draw[ring] (0,0) circle (1.8);
  \draw[ring] (0,0) circle (1.2);
  \draw[ring] (0,0) circle (0.6);
  \node[font=\scriptsize\bfseries\color{jugeoGreen!80!black}] at (0,0) {kernel};
  \node[font=\scriptsize\color{jugeoDark}] at (0,0.95) {middle layers};
  \node[font=\scriptsize\color{jugeoDark}] at (0,1.55) {user-facing};
\end{tikzpicture}
\end{center}
\end{column}
\end{columns}
\end{frame}

% ---------------------------------------------------------------------------
% Slide 165: Layer 2 — Geometry
% ---------------------------------------------------------------------------
\begin{frame}
\frametitle{Layer 2 --- Geometry}
\begin{block}{Components}
  \begin{itemize}
    \item \jugmark{Sites} --- the topological space of your codebase
    \item \jugmark{Coordinates} --- precise addresses (module.function.line)
    \item \jugmark{Covers} --- collections of inputs that span the behavior space
    \item \jugmark{Descent} --- 4 strategies for gluing local proofs into global ones
    \item \jugmark{Hypercovers} and \jugmark{supports}
  \end{itemize}
\end{block}
\vfill
\begin{center}
\begin{tikzpicture}
  \node[fill=jugeoLight, rounded corners, inner sep=8pt, text width=10cm, align=center]
    {\itshape The mathematical foundation --- the ``map'' of your codebase};
\end{tikzpicture}
\end{center}
\end{frame}

% ---------------------------------------------------------------------------
% Slide 166: Layer 3 — Judgments
% ---------------------------------------------------------------------------
\begin{frame}
\frametitle{Layer 3 --- Judgments}
\begin{columns}[T]
\begin{column}{0.55\textwidth}
\begin{block}{The 8-Tuple Algebra}
  \begin{itemize}
    \item Contexts
    \item Sections
    \item Comparisons
    \item Judgment terms
  \end{itemize}
\end{block}
\vspace{6pt}
Where \jugmark{claims about code} live.\\[4pt]
Every verification result is expressed as a judgment.
\end{column}
\begin{column}{0.42\textwidth}
\begin{center}
\begin{tikzpicture}[
  field/.style={draw=jugeoBlue, fill=jugeoLight, rounded corners,
                minimum width=28mm, minimum height=6mm,
                font=\scriptsize, text=jugeoDark, line width=0.6pt}
]
  \node[field] (f1) at (0,2.1) {context $\Gamma$};
  \node[field] (f2) at (0,1.5) {site $S$};
  \node[field] (f3) at (0,0.9) {coordinate $c$};
  \node[field] (f4) at (0,0.3) {section $\sigma$};
  \node[field] (f5) at (0,-0.3){comparison $\sim$};
  \node[field] (f6) at (0,-0.9){evidence $e$};
  \node[field] (f7) at (0,-1.5){trust $t$};
  \node[field] (f8) at (0,-2.1){mode $m$};
  \node[draw=jugeoBlue, fit=(f1)(f8), inner sep=4pt, rounded corners,
        line width=1pt, label={[font=\scriptsize\bfseries\color{jugeoBlue}]above:Judgment}] {};
\end{tikzpicture}
\end{center}
\end{column}
\end{columns}
\end{frame}

% ---------------------------------------------------------------------------
% Slide 167: Layer 4 — Evidence
% ---------------------------------------------------------------------------
\begin{frame}
\frametitle{Layer 4 --- Evidence}
\begin{block}{What It Provides}
  \begin{itemize}
    \item \jugmark{Trust algebra} --- combine, compare, and degrade trust values
    \item \jugmark{6+ evidence channels}: Z3, testing, AI, manual, runtime, static
    \item \jugmark{Certificates} --- unforgeable proofs of evidence origin
    \item \jugmark{Provenance} --- complete audit trail for every claim
    \item \jugmark{Manifests} --- packaged evidence bundles
  \end{itemize}
\end{block}
\vfill
\begin{center}
\begin{tikzpicture}
  \node[fill=jugeoLight, rounded corners, inner sep=8pt, text width=10cm, align=center]
    {\itshape The ``courtroom'' where evidence is weighed and tracked};
\end{tikzpicture}
\end{center}
\end{frame}

% ---------------------------------------------------------------------------
% Slide 168: Layer 5 — Solver
% ---------------------------------------------------------------------------
\begin{frame}
\frametitle{Layer 5 --- Solver}
\begin{columns}[T]
\begin{column}{0.55\textwidth}
\begin{block}{Z3 Integration}
  \begin{itemize}
    \item \jugmark{13 theory fragments} (arithmetic, arrays, bitvectors, \ldots)
    \item \jugmark{5 routing strategies} to pick the best fragment
    \item Session pool for parallel queries
    \item Countermodel extraction on failure
  \end{itemize}
\end{block}
\end{column}
\begin{column}{0.42\textwidth}
\begin{center}
\begin{tikzpicture}[
  frag/.style={draw=jugeoBlue, fill=jugeoLight, rounded corners,
               minimum width=18mm, minimum height=5mm,
               font=\tiny, text=jugeoDark, line width=0.5pt},
  arr/.style={-{Stealth[length=2mm]}, jugeoBlue!60, line width=0.5pt}
]
  \node[draw=jugeoOrange, fill=jugeoOrange!15, rounded corners,
        minimum width=22mm, minimum height=7mm,
        font=\small\bfseries, text=jugeoDark] (router) at (0,1.8) {Router};
  \node[frag] (f1) at (-1.2,0.4) {LIA};
  \node[frag] (f2) at (0,0.4)    {BV};
  \node[frag] (f3) at (1.2,0.4)  {Arrays};
  \node[frag] (f4) at (-0.6,-0.4){Strings};
  \node[frag] (f5) at (0.6,-0.4) {FP};
  \node[draw=jugeoGreen, fill=jugeoGreen!15, rounded corners,
        minimum width=22mm, minimum height=7mm,
        font=\small\bfseries, text=jugeoDark] (z3) at (0,-1.4) {Z3};
  \draw[arr] (router) -- (f1);
  \draw[arr] (router) -- (f2);
  \draw[arr] (router) -- (f3);
  \draw[arr] (router) -- (f4.north);
  \draw[arr] (router) -- (f5.north);
  \foreach \f in {f1,f2,f3,f4,f5} { \draw[arr] (\f) -- (z3); }
\end{tikzpicture}
\end{center}
\end{column}
\end{columns}
\end{frame}

% ---------------------------------------------------------------------------
% Slide 169: Layer 6 — Encodings
% ---------------------------------------------------------------------------
\begin{frame}
\frametitle{Layer 6 --- Encodings}
\begin{center}
  {\Large \jugmark{18 sub-packages} --- translating Python into logic}
\end{center}
\vspace{8pt}
\begin{columns}[T]
\begin{column}{0.48\textwidth}
  \begin{itemize}
    \item Deduction rules
    \item Refinement types
    \item Path conditions
    \item Heap summaries
    \item Theorem schemas
  \end{itemize}
\end{column}
\begin{column}{0.48\textwidth}
  \begin{itemize}
    \item Text encodings
    \item Tensor encodings
    \item Effect encodings
    \item Inductive encodings
    \item \ldots and 8 more
  \end{itemize}
\end{column}
\end{columns}
\vfill
\begin{center}
\begin{tikzpicture}
  \node[fill=jugeoLight, rounded corners, inner sep=8pt, text width=10cm, align=center]
    {\itshape The largest layer --- the bridge from Python semantics to Z3 formulas};
\end{tikzpicture}
\end{center}
\end{frame}

% ---------------------------------------------------------------------------
% Slide 170: Layer 7 — Python Runtime
% ---------------------------------------------------------------------------
\begin{frame}
\frametitle{Layer 7 --- Python Runtime}
\begin{block}{JuGeo's Model of Python Execution}
  \begin{itemize}
    \item \jugmark{Scope \& state} --- variable bindings, mutation tracking
    \item \jugmark{Exceptions} --- try/except/finally flow
    \item \jugmark{Async} --- coroutines, await, event loops
    \item \jugmark{Generators} --- yield, send, close
    \item \jugmark{Context managers} --- with/as protocol
    \item \jugmark{Metaobject protocol} --- \_\_getattr\_\_, descriptors, metaclasses
    \item \jugmark{Program loading} --- imports, \_\_init\_\_, sys.path
  \end{itemize}
\end{block}
\vfill
\begin{center}
  {\small\color{gray}\itshape Not a toy subset --- the \textbf{real} Python semantics}
\end{center}
\end{frame}

% ---------------------------------------------------------------------------
% Slide 171: Layer 8 — Generation
% ---------------------------------------------------------------------------
\begin{frame}
\frametitle{Layer 8 --- Generation}
\begin{columns}[T]
\begin{column}{0.55\textwidth}
\begin{block}{The Proof Search Engine}
  \begin{itemize}
    \item \jugmark{Cover design} --- choosing the right decomposition
    \item \jugmark{Construction loops} --- iterative refinement
    \item \jugmark{Inhabitant fleets} --- parallel proof candidates
    \item \jugmark{State space exploration}
    \item \jugmark{Backpressure} --- throttling runaway search
    \item \jugmark{Replay \& gluing} --- caching and composition
  \end{itemize}
\end{block}
\end{column}
\begin{column}{0.42\textwidth}
\begin{center}
\begin{tikzpicture}[
  step/.style={draw=jugeoBlue, fill=jugeoLight, rounded corners,
               minimum width=24mm, minimum height=6mm,
               font=\scriptsize\bfseries, text=jugeoDark, line width=0.6pt},
  arr/.style={-{Stealth[length=2.5mm]}, jugeoGreen!70!black, line width=0.7pt}
]
  \node[step] (s1) at (0,2.0) {Design cover};
  \node[step] (s2) at (0,1.0) {Launch fleet};
  \node[step] (s3) at (0,0.0) {Explore};
  \node[step] (s4) at (0,-1.0){Glue};
  \node[step, fill=jugeoGreen!15] (s5) at (0,-2.0){Done};
  \draw[arr] (s1)--(s2);
  \draw[arr] (s2)--(s3);
  \draw[arr] (s3)--(s4);
  \draw[arr] (s4)--(s5);
  \draw[arr, dashed, jugeoOrange] (s3.east) to[out=0,in=0] node[right,font=\tiny\color{jugeoOrange}]{retry} (s2.east);
\end{tikzpicture}
\end{center}
\end{column}
\end{columns}
\end{frame}

% ---------------------------------------------------------------------------
% Slide 172: Layer 9 — Orchestration
% ---------------------------------------------------------------------------
\begin{frame}
\frametitle{Layer 9 --- Orchestration}
\begin{columns}[T]
\begin{column}{0.55\textwidth}
\begin{block}{The Conductor}
  \begin{itemize}
    \item \jugmark{Synthesis orchestrator} --- top-level control
    \item \jugmark{Fleet management} --- allocating proof workers
    \item \jugmark{Semantic control} --- what to prove next
    \item \jugmark{Move selection} --- game-tree search over tactics
    \item \jugmark{Mixed-evidence routing}
    \item \jugmark{Frontier objectives} --- prioritizing open goals
  \end{itemize}
\end{block}
\end{column}
\begin{column}{0.42\textwidth}
\begin{center}
\begin{tikzpicture}[
  worker/.style={circle, draw=jugeoBlue, fill=jugeoLight, minimum size=7mm,
                 font=\scriptsize\bfseries, line width=0.6pt},
  arr/.style={-{Stealth[length=2mm]}, jugeoBlue!60, line width=0.6pt}
]
  \node[draw=jugeoOrange, fill=jugeoOrange!15, rounded corners,
        minimum width=24mm, minimum height=8mm,
        font=\small\bfseries, text=jugeoDark] (orch) at (0,1.5) {Orchestrator};
  \node[worker] (w1) at (-1.5,0) {W1};
  \node[worker] (w2) at (0,0)    {W2};
  \node[worker] (w3) at (1.5,0)  {W3};
  \node[worker] (w4) at (-0.75,-1.2) {W4};
  \node[worker] (w5) at (0.75,-1.2)  {W5};
  \foreach \w in {w1,w2,w3,w4,w5} { \draw[arr] (orch) -- (\w); }
\end{tikzpicture}\\[4pt]
{\scriptsize Coordinates a fleet of proof workers}
\end{center}
\end{column}
\end{columns}
\end{frame}

% ---------------------------------------------------------------------------
% Slide 173: Layer 10 — Problem Modes
% ---------------------------------------------------------------------------
\begin{frame}
\frametitle{Layer 10 --- Problem Modes}
\begin{center}
\begin{tikzpicture}[
  mode/.style={draw=jugeoBlue, fill=jugeoLight, rounded corners,
               minimum width=32mm, minimum height=7mm,
               font=\small, text=jugeoDark, line width=0.7pt}
]
  \node[mode] (m1) at (-3, 0.8) {\jugmark{Spec satisfaction}};
  \node[mode] (m2) at ( 0, 0.8) {\jugmark{Equivalence}};
  \node[mode] (m3) at ( 3, 0.8) {\jugmark{Bug detection}};
  \node[mode] (m4) at (-3,-0.4) {\jugmark{Relational refinement}};
  \node[mode] (m5) at ( 0,-0.4) {\jugmark{Repair semantics}};
  \node[mode] (m6) at ( 3,-0.4) {\jugmark{Doc generation}};
\end{tikzpicture}
\end{center}
\vfill
\begin{itemize}
  \item Each mode is a user-facing \textbf{verification goal}
  \item Same architecture underneath --- different question at the top
  \item You pick the mode; JuGeo picks the strategy
\end{itemize}
\end{frame}

% ---------------------------------------------------------------------------
% Slide 174: The CLI Interface
% ---------------------------------------------------------------------------
\begin{frame}[fragile]
\frametitle{The CLI Interface}
\begin{block}{Simple Entry Points for Complex Verification}
\begin{lstlisting}
# Prove a specification holds
$ jugeo prove myapp.auth

# Find bugs automatically
$ jugeo bugs myapp.database

# Check two implementations are equivalent
$ jugeo equiv myapp.sort_v1 myapp.sort_v2

# Encode Python into logic (inspect)
$ jugeo encode myapp.parser

# Run descent (glue local proofs)
$ jugeo descend myapp

# Load and analyze a codebase
$ jugeo load myproject/
\end{lstlisting}
\end{block}
\vfill
\begin{center}
  {\small One command per capability. Designed for the terminal.}
\end{center}
\end{frame}

% ---------------------------------------------------------------------------
% Slide 175: The Easy API
% ---------------------------------------------------------------------------
\begin{frame}[fragile]
\frametitle{The Easy API}
\begin{center}
  {\Large Seven one-function entry points}
\end{center}
\vspace{8pt}
\begin{exampleblock}{One Import, One Call}
\begin{lstlisting}
from jugeo import prove, bugs, equiv
from jugeo import ideate, carry, spec, research

# Prove a function meets its spec
result = prove("myapp.auth.validate_token")

# Find bugs in a module
issues = bugs("myapp.database")

# Check equivalence of two implementations
ok = equiv("myapp.sort_v1", "myapp.sort_v2")
\end{lstlisting}
\end{exampleblock}
\vfill
\begin{center}
  \jugmark{One function per capability.} Designed for Python developers.
\end{center}
\end{frame}

% ---------------------------------------------------------------------------
% Slide 176: Experiments and Benchmarks
% ---------------------------------------------------------------------------
\begin{frame}
\frametitle{Experiments and Benchmarks}
\begin{center}
\begin{tikzpicture}[
  bignum/.style={font=\fontsize{42}{50}\selectfont\bfseries\color{jugeoBlue}}
]
  \node[bignum] (n) {70+};
  \node[font=\large\color{jugeoDark}, below=2pt of n] {experiments};
\end{tikzpicture}
\end{center}
\vspace{4pt}
\begin{columns}[T]
\begin{column}{0.48\textwidth}
  \begin{itemize}
    \item Site scaling
    \item Judgment algebra
    \item Descent obstructions
    \item Cover design
  \end{itemize}
\end{column}
\begin{column}{0.48\textwidth}
  \begin{itemize}
    \item Bug detection
    \item Heap aliasing
    \item Countermodel extraction
    \item \ldots and many more
  \end{itemize}
\end{column}
\end{columns}
\vfill
\begin{center}
  {\small\color{gray} All in \texttt{experiments/} --- reproducible with one command}
\end{center}
\end{frame}

% ---------------------------------------------------------------------------
% Slide 177: Mathematical Foundations
% ---------------------------------------------------------------------------
\begin{frame}
\frametitle{Mathematical Foundations}
\begin{columns}[T]
\begin{column}{0.55\textwidth}
\begin{block}{Formal Theory}
  \begin{itemize}
    \item \jugmark{67-chapter} formal theory document
    \item Covers: sites, judgments, descent, trust algebra, cut elimination, effects
    \item Written in \LaTeX\ (\texttt{theory2.tex})
  \end{itemize}
\end{block}
\vspace{6pt}
\begin{exampleblock}{Lean 4 Formalizations}
  Machine-checked proofs in \texttt{proofs/lean/}\\[4pt]
  Core theorems verified independently.
\end{exampleblock}
\end{column}
\begin{column}{0.42\textwidth}
\begin{center}
\begin{tikzpicture}[
  doc/.style={draw=jugeoBlue, fill=jugeoLight, rounded corners,
              minimum width=26mm, minimum height=8mm,
              font=\small\bfseries, text=jugeoDark, line width=0.6pt},
  arr/.style={-{Stealth[length=2.5mm]}, jugeoBlue!60, line width=0.6pt}
]
  \node[doc] (th) at (0,1.5) {theory2.tex};
  \node[doc] (ln) at (0,0)   {Lean 4 proofs};
  \node[doc, fill=jugeoGreen!15] (impl) at (0,-1.5) {Python impl};
  \draw[arr] (th) -- node[right, font=\tiny\color{gray}]{formalizes} (ln);
  \draw[arr] (ln) -- node[right, font=\tiny\color{gray}]{validates} (impl);
\end{tikzpicture}
\end{center}
\end{column}
\end{columns}
\end{frame}

% ---------------------------------------------------------------------------
% Slide 178: The Web Interface
% ---------------------------------------------------------------------------
\begin{frame}
\frametitle{The Web Interface}
\begin{center}
\begin{tikzpicture}[
  page/.style={draw=jugeoBlue, fill=jugeoLight, rounded corners,
               minimum width=30mm, minimum height=9mm,
               font=\small\bfseries, text=jugeoDark, line width=0.7pt},
  arr/.style={-{Stealth[length=2.5mm]}, jugeoBlue!50, line width=0.6pt}
]
  \node[draw=jugeoDark, fill=jugeoDark!10, rounded corners,
        minimum width=34mm, minimum height=10mm,
        font=\large\bfseries, text=jugeoDark] (hub) at (0,2) {web/};
  \node[page] (tut) at (-3.2,0) {Tutorials};
  \node[page] (api) at (-1.1,0) {API Reference};
  \node[page] (con) at (1.1,0)  {Concept Guides};
  \node[page] (cmp) at (3.2,0)  {Comparisons};
  \draw[arr] (hub) -- (tut);
  \draw[arr] (hub) -- (api);
  \draw[arr] (hub) -- (con);
  \draw[arr] (hub) -- (cmp);
\end{tikzpicture}
\end{center}
\vfill
\begin{itemize}
  \item Full documentation site in the project's \texttt{web/} directory
  \item Step-by-step tutorials for new users
  \item API reference auto-generated from docstrings
  \item Comparison pages: JuGeo vs.\ LEAN, Coq, F*, Dafny
\end{itemize}
\end{frame}

% ---------------------------------------------------------------------------
% Slide 179: Testing Infrastructure
% ---------------------------------------------------------------------------
\begin{frame}
\frametitle{Testing Infrastructure}
\begin{columns}[T]
\begin{column}{0.55\textwidth}
\begin{block}{Scale}
  \begin{itemize}
    \item \jugmark{300+ benchmark cases} across 8 families
    \item Integration tests for CLI robustness
    \item Directed research tests
    \item Research orchestration tests
    \item Comprehensive \texttt{pytest} suite
  \end{itemize}
\end{block}
\end{column}
\begin{column}{0.42\textwidth}
\begin{center}
\begin{tikzpicture}[
  fam/.style={draw=jugeoGreen, fill=jugeoGreen!10, rounded corners,
              minimum width=24mm, minimum height=5.5mm,
              font=\scriptsize, text=jugeoDark, line width=0.5pt}
]
  \node[fam] (f1) at (0,2.45) {Unit tests};
  \node[fam] (f2) at (0,1.75) {Integration tests};
  \node[fam] (f3) at (0,1.05) {CLI tests};
  \node[fam] (f4) at (0,0.35) {Benchmark cases};
  \node[fam] (f5) at (0,-0.35){Research tests};
  \node[fam] (f6) at (0,-1.05){Orchestration};
  \node[fam] (f7) at (0,-1.75){Regression};
  \node[fam] (f8) at (0,-2.45){End-to-end};
  \node[draw=jugeoGreen, fit=(f1)(f8), inner sep=4pt, rounded corners,
        line width=1pt, label={[font=\scriptsize\bfseries\color{jugeoGreen}]above:8 Families}] {};
\end{tikzpicture}
\end{center}
\end{column}
\end{columns}
\end{frame}

% ---------------------------------------------------------------------------
% Slide 180: Deduction Rules — First-Class Inference
% ---------------------------------------------------------------------------
\begin{frame}
\frametitle{Deduction Rules --- First-Class Inference}
\begin{columns}[T]
\begin{column}{0.48\textwidth}
\begin{alertblock}{Structural Rules}
  \begin{itemize}
    \item Weakening
    \item Contraction
    \item Exchange
    \item Cut
  \end{itemize}
\end{alertblock}
\vspace{4pt}
\begin{exampleblock}{Semantic Rules}
  \begin{itemize}
    \item Introduction
    \item Elimination
    \item Computation
  \end{itemize}
\end{exampleblock}
\end{column}
\begin{column}{0.48\textwidth}
\begin{block}{Judgment Transitions}
  Rules that move between judgment forms:
  \begin{itemize}
    \item Verification $\to$ synthesis
    \item Local $\to$ global (descent)
    \item Single $\to$ relational
  \end{itemize}
\end{block}
\vspace{4pt}
\begin{center}
\begin{tikzpicture}
  \node[fill=jugeoLight, rounded corners, inner sep=6pt, font=\footnotesize,
        text width=38mm, align=center]
    {All rules are \jugmark{first-class Python objects} --- inspect, compose, and replay them};
\end{tikzpicture}
\end{center}
\end{column}
\end{columns}
\end{frame}
